\documentclass[12pt]{article}

% Keep the template's layout with fonts available in the local TeX installation.
\usepackage[T1]{fontenc}
\usepackage{times}
\usepackage{microtype}
\usepackage{amsmath}
\usepackage[parfill]{parskip}
\setlength{\parindent}{0pt}
\setlength{\parskip}{2ex plus 0.2ex minus 0.2ex}
\usepackage[letterpaper,margin=1in]{geometry}
\usepackage[hidelinks]{hyperref}

\begin{document}

\begin{flushleft}
\textbf{Week 2: Linear and Logistic Regression} \\
Jaisel Singh (js6897) \\
September 23, 2026
\end{flushleft}

\vspace{0.1in}

What stood out to me in \emph{Regression and Other Stories} was the distinction
between using regression for prediction and using it to compare groups. A
coefficient can describe how average earnings differ between taller and shorter
people, but that alone does not tell us what would happen to someone's earnings
if their height changed. Calling a coefficient an ``effect'' can therefore hide
assumptions that the fitted model does not establish. This makes interpreting a
regression seem at least as important as fitting it.

I also found it helpful to think of regression as modeling a distribution of
outcomes given the inputs. With independent normal errors of equal variance,
$y_i \mid x_i \sim \mathcal{N}(a+bx_i,\sigma^2)$, maximizing the likelihood over
$a$ and $b$ gives the same estimates as minimizing squared residuals. This
connects a familiar fitting rule to an explicit assumption about the data,
although normality is not required simply to fit least squares. Keeping the
distribution in mind also clarified regression to the mean: predicted heights
can be closer to the average without implying that everyone eventually has the
same height, because actual outcomes still vary around those predictions.

Logistic regression extends this perspective to binary outcomes. It models a
Bernoulli probability $p(x)$, with $\log[p(x)/(1-p(x))]=a+bx$, so predicted
probabilities stay between zero and one. The coefficient describes a constant
difference in log-odds per unit of $x$; the corresponding difference in
probability depends on the starting value of $x$. I find the shared likelihood
view useful: choosing a distribution for the outcome determines how we fit the
model and interpret its predictions. The question I am left with is how to tell
when a model predicts reasonably well but gives misleading comparisons because
important variables or interactions have been left out.

\vspace{0.05in}
{\footnotesize
\textit{Readings:} Gelman, Hill, and Vehtari, \href{https://avehtari.github.io/ROS-Examples/}{\emph{Regression and Other Stories}};
Shalizi, \href{https://www.stat.cmu.edu/~cshalizi/ADAfaEPoV/}{\emph{Advanced Data Analysis from an Elementary Point of View}} (logistic regression).
\par}

\end{document}

%%% Local Variables:
%%% mode: latex
%%% TeX-master: t
%%% End:
